\section{Related Work}
\label{sec:related_work}

\subsection{Reinforcement Learning from Human Feedback}

RLHF was introduced for language model alignment by Ziegler et al. \citep{ziegler2019fine} and scaled to production systems with InstructGPT \citep{ouyang2022instructgpt}. The standard pipeline involves: (1) supervised fine-tuning on demonstration data, (2) training a reward model on human preference rankings, and (3) optimizing the policy via PPO against the reward model with KL regularization.

Constitutional AI (CAI) \citep{bai2022constitutional} introduced an alternative: using AI feedback (RLAIF) with explicit constitutional principles for harmlessness training. CAI's key insight---that explicit rules can induce implicit behavioral improvements---motivates our hypothesis that explicit constraint training transfers to implicit safety.

Direct Preference Optimization (DPO) \citep{rafailov2023dpo} eliminates the reward model, optimizing directly on preference pairs. While DPO simplifies training, it remains single-objective and does not address multi-signal integration.

\subsection{Multi-Objective Alignment}

Askell et al. \citep{askell2021anthropic} formalized the HHH framework (Helpful, Harmless, Honest), analyzing trade-offs between objectives. Their work demonstrated that multi-objective optimization is feasible but focused on helpfulness-harmlessness, not helpfulness-controllability.

Gao et al. \citep{gao2022scaling} characterized reward model overoptimization, showing that excessive optimization degrades true reward. This informs our use of KL regularization.

\subsection{Instruction Following and Controllability}

IFEval \citep{zhou2023ifeval} introduced a benchmark for verifiable instruction following, measuring compliance with 25 constraint types (format, length, keywords, structure). We extend IFEval's use from evaluation to training signal extraction.

AlpacaEval \citep{dubois2024alpacaeval} measures instruction-following quality via GPT-4 judgment against reference outputs. We use AlpacaEval LC as our helpfulness metric.

\subsection{Bidirectional Alignment Framework}

Sun et al. \citep{sun2024bidirectional} survey 400+ papers to propose a bidirectional taxonomy:
\begin{itemize}
    \item \textbf{AI$\to$Human:} Specification integration---models adopting human values
    \item \textbf{Human$\to$AI:} Agency preservation---humans retaining control
\end{itemize}

Their framework is theoretical; we provide the first empirical implementation.
